\documentclass[11pt,a4paper]{scrartcl}

\usepackage{amssymb}
\usepackage{amsmath}
\usepackage{yhmath}
\usepackage{tikz}
\usetikzlibrary{angles,quotes,calc,decorations.markings,intersections}

\usepackage[utf8]{inputenc}
\usepackage[T1]{fontenc}
\newcommand{\ats}[1]{\par\noindent\textit{Atsakymas:} #1}
\usepackage{cancel}
\usepackage{tikz}
\usepackage{lmodern} 
\usepackage{amsmath,amssymb,amsthm}
\usepackage{graphicx}
\usepackage[dvipsnames,svgnames]{xcolor}
\usepackage{hyperref}
\usepackage{mathrsfs}
\usepackage[shortlabels]{enumitem}
\usepackage[obeyFinal,textsize=scriptsize,shadow,loadshadowlibrary]{todonotes}
\usepackage{mathtools}
\usepackage{microtype}

%%% Paragraph layout
\setlength{\parindent}{0pt}
\setlength{\parskip}{0.6\baselineskip}

%%% Colored sections
\renewcommand*{\sectionformat}%
  {\color{purple}\S\thesection\enskip}
\renewcommand*{\subsectionformat}%
  {\color{purple}\S\thesubsection\enskip}
\renewcommand*{\subsubsectionformat}%
  {\color{purple}\S\thesubsubsection\enskip}
\addtokomafont{paragraph}{\color{orange!35!black}\P\ }
\KOMAoptions{numbers=noenddot}

%%% Title
\addtokomafont{subtitle}{\Large}
\setkomafont{author}{\Large\scshape}
\setkomafont{date}{\Large\normalsize}
% Neigiamas vertikalus tarpas, kuris pakelia dokumento antraštę į viršų. 
% Galite keisti -2.5cm į -3cm (aukščiau) arba -1.5cm (žemiau).
\titlehead{\vspace*{-7.5cm}} 

% PRIDĖTA: Automatiškai perrašome datą į mokyklos pavadinimą
\AtBeginDocument{%
  \date{\normalsize Vilniaus licėjus}
}

%%% Page setup
\usepackage[headsepline]{scrlayer-scrpage}
\renewcommand{\headfont}{}
\addtolength{\textheight}{3.14cm}
\setlength{\footskip}{0.5in}
\setlength{\headsep}{10pt}
% Puslapio antraštė turi rodyti tik konkrečius dokumento duomenis.
% Nustatome juos aiškiai, kad neatsirastų "\@author" / "\@title" arba senų reikšmių.
\makeatletter
\AtBeginDocument{%
  \ihead{\footnotesize\textbf{Andrius Berniukevičius}}%
  \ohead{\footnotesize\textbf{\@title}}%
}
\makeatother
\automark{section}
\chead{}
\cfoot{\pagemark}

%%% Macros
\providecommand{\ol}{\overline}
\providecommand{\ul}{\underline}
\providecommand{\wt}{\widetilde}
\providecommand{\wh}{\widehat}
\providecommand{\eps}{\varepsilon}
\providecommand{\half}{\frac{1}{2}}
\providecommand{\inv}{^{-1}}
\newcommand{\dang}{\measuredangle} %% Directed angle
\newcommand{\vocab}[1]{\textbf{\color{blue}\sffamily #1}}
\providecommand{\CC}{\mathbb C}
\providecommand{\FF}{\mathbb F}
\providecommand{\NN}{\mathbb N}
\providecommand{\QQ}{\mathbb Q}
\providecommand{\RR}{\mathbb R}
\providecommand{\ZZ}{\mathbb Z}
\providecommand{\ts}{\textsuperscript}
\providecommand{\dg}{^\circ}
\providecommand{\ii}{\item}
\providecommand{\defeq}{\coloneqq}
\DeclareMathOperator*{\lcm}{lcm}
\DeclareMathOperator*{\argmin}{arg min}
\DeclareMathOperator*{\argmax}{arg max}
\providecommand{\hrulebar}{\par
\hspace{\fill}\rule{0.95\linewidth}{.7pt}\hspace{\fill}
\par\nointerlineskip \vspace{\baselineskip}}

%%% Optional colored theorems
\usepackage{thmtools}
\usepackage[framemethod=TikZ]{mdframed}
\mdfdefinestyle{mdbluebox}{%
  roundcorner=10pt,
  linewidth=1pt,
  skipabove=12pt,
  innerbottommargin=9pt,
  skipbelow=2pt,
  linecolor=blue,
  nobreak=true,
  backgroundcolor=TealBlue!5,
}
\declaretheoremstyle[
  headfont=\sffamily\bfseries\color{MidnightBlue},
  mdframed={style=mdbluebox},
  headpunct={\\[3pt]},
  postheadspace={0pt}
]{thmbluebox}
\mdfdefinestyle{mdredbox}{%
  linewidth=0.5pt,
  skipabove=12pt,
  frametitleaboveskip=5pt,
  frametitlebelowskip=0pt,
  skipbelow=2pt,
  frametitlefont=\bfseries,
  innertopmargin=4pt,
  innerbottommargin=8pt,
  nobreak=true,
  backgroundcolor=Salmon!5,
  linecolor=RawSienna,
}
\declaretheoremstyle[
  headfont=\bfseries\color{RawSienna},
  mdframed={style=mdredbox},
  headpunct={\\[3pt]},
  postheadspace={0pt},
]{thmredbox}
\mdfdefinestyle{mdgreenbox}{%
  skipabove=8pt,
  linewidth=2pt,
  rightline=false,
  leftline=true,
  topline=false,
  bottomline=false,
  linecolor=ForestGreen,
  backgroundcolor=ForestGreen!5,
}
\declaretheoremstyle[
  headfont=\bfseries\sffamily\color{ForestGreen!70!black},
  bodyfont=\normalfont,
  spaceabove=2pt,
  spacebelow=1pt,
  mdframed={style=mdgreenbox},
  headpunct={ --- },
]{thmgreenbox}
\mdfdefinestyle{mdblackbox}{%
  skipabove=8pt,
  linewidth=3pt,
  rightline=false,
  leftline=true,
  topline=false,
  bottomline=false,
  linecolor=black,
  backgroundcolor=RedViolet!5!gray!5,
}
\declaretheoremstyle[
  headfont=\bfseries,
  bodyfont=\normalfont\small,
  spaceabove=0pt,
  spacebelow=0pt,
  mdframed={style=mdblackbox}
]{thmblackbox}

%% Aplinkos su rėmeliais (thmtools)
\declaretheorem[style=thmbluebox,name=Teorema]{theorem}
\declaretheorem[style=thmbluebox,name=Lema]{lemma}
\declaretheorem[style=thmbluebox,name=Savybė]{property}
\declaretheorem[style=thmbluebox,name=Teiginys]{proposition}
\declaretheorem[style=thmbluebox,name=Išvada]{corollary}
\declaretheorem[style=thmbluebox,name=Teorema,numbered=no]{theorem*}
\declaretheorem[style=thmbluebox,name=Lema,numbered=no]{lemma*}
\declaretheorem[style=thmbluebox,name=Teiginys,numbered=no]{proposition*}
\declaretheorem[style=thmbluebox,name=Išvada,numbered=no]{corollary*}
\declaretheorem[style=thmgreenbox,name=Algoritmas]{algorithm}
\declaretheorem[style=thmgreenbox,name=Algoritmas,numbered=no]{algorithm*}
\declaretheorem[style=thmgreenbox,name=Teiginys]{claim}
\declaretheorem[style=thmgreenbox,name=Teiginys,numbered=no]{claim*}
\declaretheorem[style=thmredbox,name=Pavyzdys]{example}
\declaretheorem[style=thmredbox,name=Pavyzdys,numbered=no]{example*}
\declaretheorem[style=thmblackbox,name=Pastaba]{remark}
\declaretheorem[style=thmblackbox,name=Pastaba,numbered=no]{remark*}

%% Standartinės amsthm aplinkos (Apibrėžimai ir užduotys)
\theoremstyle{definition}
\ifcsname conjecture\endcsname\else\newtheorem{conjecture}{Hipotezė}\fi
\ifcsname definition\endcsname\else\newtheorem{definition}{Apibrėžimas}\fi
\ifcsname fact\endcsname\else\newtheorem{fact}{Faktas}\fi
\ifcsname answer\endcsname\else\newtheorem{answer}{Atsakymas}\fi
\ifcsname ques\endcsname\else\newtheorem{ques}{Klausimas}\fi
\ifcsname exercise\endcsname\else\newtheorem{exercise}{Užduotis}\fi
\ifcsname problem\endcsname\else\newtheorem{problem}{Uždavinys}[section]\fi
\ifcsname conjecture*\endcsname\else\newtheorem*{conjecture*}{Hipotezė}\fi
\ifcsname definition*\endcsname\else\newtheorem*{definition*}{Apibrėžimas}\fi
\ifcsname fact*\endcsname\else\newtheorem*{fact*}{Faktas}\fi
\ifcsname answer*\endcsname\else\newtheorem*{answer*}{Atsakymas}\fi
\ifcsname ques*\endcsname\else\newtheorem*{ques*}{Klausimas}\fi
\ifcsname exercise*\endcsname\else\newtheorem*{exercise*}{Užduotis}\fi
\ifcsname problem*\endcsname\else\newtheorem*{problem*}{Uždavinys}\fi

\newcounter{answerssection}
\newcommand{\answerssection}{%
  \setcounter{answerssection}{\value{section}}%
  \section{Atsakymai}%
  \setcounter{problem}{0}%
  \renewcommand{\theproblem}{\arabic{answerssection}.\arabic{problem}}%
}

%% GALUTINIS SPRENDIMAS PROOF APLINKAI
%% --- PRADŽIA: Įrodymo aplinkos sutvarkymas ---
\makeatletter

% 1. Užtikriname, kad žodis būtų "Įrodymas", net jei "babel" paketas bando jį perrašyti
\AtBeginDocument{%
  \renewcommand{\proofname}{Įrodymas}%
  \@ifpackageloaded{babel}{%
    \addto\captionslithuanian{\renewcommand{\proofname}{Įrodymas}}%
  }{}%
}

% 2. Priverstinai pašaliname scrartcl klasės bold šriftą ir paliekame TIK italic
% 3. Sujungiame su tuščiu kvadratėliu dešiniajame kampe.
\renewcommand{\qedsymbol}{\ensuremath{\Box}}
\renewenvironment{proof}[1][\proofname]{\par
  \pushQED{\qed}%
  \normalfont \topsep6\p@\@plus6\p@\relax
  \trivlist
  \item[\hskip\labelsep
        \normalfont\itshape % Priverčia naudoti tik pasvirą šriftą, kaip nuotraukoje
    #1\@addpunct{.}]\ignorespaces
}{%
  \popQED\endtrivlist\@endpefalse
}
\newenvironment{solution}{\par
  \normalfont \topsep6\p@\@plus6\p@\relax
  \trivlist
  \item[\hskip\labelsep
        \normalfont\itshape
    \textit{Sprendimas}\@addpunct{.}]\ignorespaces
}{%
  \endtrivlist\@endpefalse
}
\newcommand{\ans}[1]{\par\noindent\textit{Atsakymas:} #1}
\makeatother


\title{Geometrija: Kampų medžioklė apskritime}
\author{Andrius Berniukevičius}

\newenvironment{solution}
    {\begin{list}{}{\setlength{\leftmargin}{10pt}\setlength{\rightmargin}{10pt}}\item[]}
    {\end{list}}

\begin{document}

\maketitle

\section{Naujas lygis: Nematomas ryšys}

Pirmajame užsiėmime išmokome sekti kampus $\alpha$ per lygiašonius trikampius ir lygiagrečias tieses. Tačiau ten kampą galėjome perkelti tik į gretimą, besiliečiančią figūrą. 

Olimpiadinėje geometrijoje egzistuoja įrankis, leidžiantis paimti kampą brėžinio viršuje ir akimirksniu, be jokių tarpinių skaičiavimų, perkelti jį į patį brėžinio dugną. Šis įrankis – \textbf{apskritimas}. Apskritimas veikia kaip idealus veidrodis, atspindintis kampus į skirtingas erdves. Net jei uždavinio sąlygoje apskritimo nėra, mes dažnai jį nusibrėžiame patys!

% ==========================================
% BAZINIAI GINKLAI 2
% ==========================================
\section{Apskritimo anatomija: 3 nauji mechanizmai}

Pamirškite $C = 2\pi R$ ar ploto formules. Kampų medžioklėje mums rūpi tik viena apskritimo savybė – kaip jis keičia ir perkelia kampus. Prieš formuluojant taisykles, įsiveskime tris esmines sąvokas.

\subsection*{Bazinės sąvokos}

\begin{definition}[Centrinis kampas]
Kampas, kurio viršūnė yra apskritimo centras, o kampo kraštinės yra apskritimo spinduliai.
\end{definition}

\begin{center}
\begin{tikzpicture}[scale=1.2]
    \coordinate (O) at (0,0);
    \coordinate (A) at (210:2);
    \coordinate (B) at (330:2);
    \draw[thick, dashed] (O) circle (2);
    \draw[thick, blue] (A) -- (O) -- (B);
    \pic [draw, fill=blue!10, fill opacity=0.35, angle radius=8mm,
        pic text options={text=black}] {angle = A--O--B};
    \filldraw (O) circle (1.2pt) node[above] {$O$};
    \filldraw (A) circle (1.2pt) node[below left] {$A$};
    \filldraw (B) circle (1.2pt) node[below right] {$B$};
\end{tikzpicture}
\end{center}

\begin{definition}[Lanko dydis ir jo žymėjimas]
Lanko laipsninis matas yra visada lygus jį atitinkančiam centriniam kampui. Lankas žymimas lenktu ženklu virš jo galų. Jei centrinis kampas $\angle AOB = 120^\circ$, tai į jį besiremiančio lanko $AB$ dydis rašomas taip:
\[ \wideparen{AB} = 120^\circ \]
\end{definition}

\begin{center}
\begin{tikzpicture}[scale=1.2]
    \coordinate (O) at (0,0);
    \coordinate (A) at (210:2);
    \coordinate (B) at (330:2);
    \draw[thick, dashed] (O) circle (2);
    \draw[ultra thick, red!85!black] (210:2) arc (210:330:2);
    \draw[very thick, blue] (A) -- (O) -- (B);
    \pic [draw=blue, fill=blue!30, fill opacity=0.75, "$\mathbf{120^\circ}$", angle radius=8mm,
        angle eccentricity=1.35, pic text options={text=black, font=\bfseries}] {angle = A--O--B};
    \filldraw (A) circle (1.2pt) node[below left] {$A$};
    \filldraw (B) circle (1.2pt) node[below right] {$B$};
    \node[red, below, yshift=-1mm] at (270:2) {$\wideparen{AB} = 120^\circ$};
\end{tikzpicture}
\end{center}

\begin{definition}[Įbrėžtinis kampas]
Kampas, kurio viršūnė yra ant apskritimo, o kampo kraštinės yra apskritimo stygos.
\end{definition}

\begin{center}
\begin{tikzpicture}[scale=1.2]
    \coordinate (A) at (210:2);
    \coordinate (B) at (330:2);
    \coordinate (C) at (70:2);
    \draw[thick, dashed] (0,0) circle (2);
    \draw[thick, green!60!black] (A) -- (C) -- (B);
    \pic [draw, fill=green!10, fill opacity=0.35, angle radius=10mm,
        pic text options={text=black}] {angle = A--C--B};
    \filldraw (A) circle (1.2pt) node[below left] {$A$};
    \filldraw (B) circle (1.2pt) node[below right] {$B$};
    \filldraw (C) circle (1.2pt) node[above] {$C$};
\end{tikzpicture}
\end{center}
\vspace{0.3cm}

\subsection*{1. Centrinio - įbrėžtinio kampų savybė}
Jei abu kampai remiasi į tą patį lanką, centrinis kampas visada yra lygiai du kartus didesnis už įbrėžtinį.
\begin{center}
\begin{tikzpicture}[scale=1.2]
    \coordinate (O) at (0,0);
    \coordinate (A) at (210:2);
    \coordinate (B) at (330:2);
    \coordinate (C) at (90:2);
    
    \draw[thick, dashed] (O) circle (2);
    \draw[thick] (A) -- (O) -- (B);
    \draw[thick] (A) -- (C) -- (B);
    
    \draw[ultra thick, red] (210:2) arc (210:330:2);
    \foreach \p in {O,A,B,C} {\filldraw (\p) circle (1.2pt);}
    
    \pic [draw, fill=blue!10, fill opacity=0.35, "$\alpha$", angle radius=8mm,
           pic text options={text=black,text opacity=1}] {angle = A--C--B};
    \pic [draw, fill=red!10, fill opacity=0.35, "$2\alpha$", angle radius=6mm,
           pic text options={text=black,text opacity=1}] {angle = A--O--B};
    
    \filldraw (O) circle (1.5pt) node[above] {$O$};
    \node[below left] at (A) {$A$};
    \node[below right] at (B) {$B$};
    \node[above] at (C) {$C$};
\end{tikzpicture}
\end{center}

Kadangi abu kampai remiasi į lanką $\wideparen{AB}$, galime sakyti ir lanko dydžio kalba: įbrėžtinio kampo dydis yra pusė jį atitinkančio lanko dydžio, o centrinio kampo dydis lygus visam lankui:
\[
\angle ACB = \frac{1}{2}\,\wideparen{AB},
\qquad
\angle AOB = \wideparen{AB}.
\]

\subsection*{2. Įbrėžtinių kampų savybė}
Tai pati dažniausiai naudojama olimpiadinė savybė. Visi įbrėžtiniai kampai, kurie „žiūri“ į tą patį lanką (arba remiasi į tą pačią stygą iš tos pačios pusės), yra lygūs. Įsivaizduokite lanką kaip kino teatro ekraną – nesvarbu, kurioje apskritimo vietoje atsisėsite, ekraną matysite tokiu pačiu kampu $\alpha$.
\begin{center}
\begin{tikzpicture}[scale=1.5]
    \draw[thick] (0,0) circle (2);
    
    \coordinate (A) at (220:2);
    \coordinate (B) at (320:2);
    \coordinate (C1) at (50:2);
    \coordinate (C2) at (100:2);
    \coordinate (C3) at (150:2);
    
    \draw[ultra thick, red] (220:2) arc (220:320:2);
    \draw[thick] (A) -- (B);
    
    \draw[thick, gray] (A) -- (C1) -- (B);
    \draw[thick, gray] (A) -- (C2) -- (B);
    \draw[thick, gray] (A) -- (C3) -- (B);
    \foreach \p in {O,A,B,C1,C2,C3} {\filldraw (\p) circle (1.2pt);}
    
    \pic [draw, fill=blue!10, fill opacity=0.4, "$\alpha$", angle radius=8mm, pic text options={text=black,text opacity=1}] {angle = A--C1--B};
    \pic [draw, fill=blue!10, fill opacity=0.4, "$\alpha$", angle radius=8mm, pic text options={text=black,text opacity=1}] {angle = A--C2--B};
    \pic [draw, fill=blue!10, fill opacity=0.4, "$\alpha$", angle radius=8mm, pic text options={text=black,text opacity=1}] {angle = A--C3--B};
    
    \node[below left] at (A) {$A$};
    \node[below right] at (B) {$B$};
\end{tikzpicture}
\end{center}

\subsection*{3. Įbrėžtinio kampo, besiremiančio į skersmenį, savybė}
Jeigu įbrėžtinis kampas remiasi į apskritimo skersmenį, tai tas kampas \textbf{visada} yra status.
\begin{center}
\begin{tikzpicture}[scale=1.5]
    \coordinate (O) at (0,0);
    \draw[thick] (O) circle (2);
    
    \coordinate (A) at (180:2);
    \coordinate (B) at (0:2);
    \coordinate (C) at (60:2);
    
    \draw[thick] (A) -- (B);
    \draw[thick] (A) -- (C) -- (B);
    \foreach \p in {O,A,B,C} {\filldraw (\p) circle (1.2pt);}
    
    % Statusis kampas (kvadračiukas)
    \coordinate (CdirA) at ($(C)!4mm!(A)$);
    \coordinate (CdirB) at ($(C)!4mm!(B)$);
    \coordinate (Crect) at ($(CdirA) + (CdirB) - (C)$);
    \draw[thick, fill=red!20, fill opacity=0.35] (CdirA) -- (Crect) -- (CdirB) -- (C) -- cycle;
    
    \filldraw (O) circle (1.5pt) node[below] {$O$};
    \node[left] at (A) {$A$};
    \node[right] at (B) {$B$};
    \node[above right] at (C) {$C$};
\end{tikzpicture}
\end{center}

% ==========================================
% PAVYZDŽIAI
% ==========================================
\section{Klasika: Kampų medžioklė apskritime}

Atidžiai išnagrinėkime, kaip šios savybės taikomos sprendžiant uždavinius. 

\begin{example}[1. Centrinio ir įbrėžtinio kampų savybės įrodymas]
Įrodykite, kad centrinis kampas visada yra dvigubai didesnis už įbrėžtinį, besiremiantį į tą patį lanką. (Sujungsime Geometrija I ir Geometrija II žinias!).
\end{example}

\textbf{Sprendimas:}\\
Šis įrodymas nereikalauja jokių naujų formulių – tik praeito užsiėmimo ginklų! 

\textbf{1 Žingsnis (Pasiruošimas):} Nusibraižykime apskritimą su centru $O$. Pažymėkime įbrėžtinį kampą $\angle ACB$ ir centrinį kampą $\angle AOB$. Per viršūnę $C$ ir centrą $O$ nubrėžkime pagalbinę liniją (skersmenį) $CD$. Ji padalina mūsų įbrėžtinį kampą į dvi dalis: $\alpha$ ir $\beta$.

\begin{center}
\begin{tikzpicture}[scale=1.2]
    \coordinate (O) at (0,0);
    \draw[thick] (O) circle (2);
    
    \coordinate (A) at (210:2);
    \coordinate (B) at (330:2);
    \coordinate (C) at (70:2);
    \coordinate (D) at (250:2); % skersmens pabaiga
    
    \draw[thick, dashed] (C) -- (D);
    \draw[thick] (A) -- (C) -- (B);
    \draw[thick] (A) -- (O) -- (B);
    \draw[thick, blue] (O) -- (C);
        \foreach \p in {O,A,B,C} {\filldraw (\p) circle (1.2pt);}
    
    \pic [draw, fill=blue!20, fill opacity=0.35, "$\alpha$", angle radius=10mm, angle eccentricity=1.35,
        pic text options={text=black,text opacity=1}] {angle = A--C--D};
    \pic [draw, fill=yellow!25, fill opacity=0.35, "$\beta$", angle radius=9mm, angle eccentricity=1.35,
        pic text options={text=black,text opacity=1}] {angle = D--C--B};
    
    \filldraw (O) circle (1pt) node[above right] {$O$};
    \node[below left] at (A) {$A$};
    \node[below right] at (B) {$B$};
    \node[above] at (C) {$C$};
    \node[below] at (D) {$D$};
\end{tikzpicture}
\end{center}

\textbf{2 Žingsnis (Klonavimas):} Pastebėkime, kad atkarpos $OA, OB$ ir $OC$ yra to paties apskritimo spinduliai! Vadinasi, $OA = OC$ ir $OB = OC$. 
Tai sukuria du lygiašonius trikampius: $\triangle AOC$ ir $\triangle BOC$.
Jei $\triangle AOC$ lygiašonis, tai jo kampai prie pagrindo lygūs: $\angle OAC = \angle OCA = \mathbf{\alpha}$.
Lygiai taip pat, $\triangle BOC$ kampai prie pagrindo: $\angle OBC = \angle OCB = \mathbf{\beta}$.

\begin{center}
\begin{tikzpicture}[scale=1.2]
    \coordinate (O) at (0,0);
    \draw[thick] (O) circle (2);
    \coordinate (A) at (210:2); \coordinate (B) at (330:2); \coordinate (C) at (70:2); \coordinate (D) at (250:2);
    \draw[thick, dashed] (C) -- (D); \draw[thick] (A) -- (C) -- (B); \draw[thick] (A) -- (O) -- (B); \draw[thick, blue] (O) -- (C);
    \foreach \p in {O,A,B,C,D} {\filldraw (\p) circle (1.2pt);}
    
    \draw[thick, blue] ($(O)!0.5!(A) + (-0.05, 0.08)$) -- ($(O)!0.5!(A) + (0.05, -0.08)$);
    \draw[thick, blue] ($(O)!0.5!(C) + (-0.085, 0.031)$) -- ($(O)!0.5!(C) + (0.085, -0.031)$);
    \draw[thick, blue] ($(O)!0.5!(B) + (0.05, 0.08)$) -- ($(O)!0.5!(B) + (-0.05, -0.08)$);
    
    \pic [draw, fill=blue!20, fill opacity=0.35, "$\alpha$", angle radius=10mm, angle eccentricity=1.35,
        pic text options={text=black,text opacity=1}] {angle = A--C--D};
    \pic [draw, fill=blue!20, fill opacity=0.35, "$\alpha$", angle radius=8mm, angle eccentricity=1.35,
        pic text options={text=black,text opacity=1}] {angle = O--A--C};
    \pic [draw, fill=yellow!25, fill opacity=0.35, "$\beta$", angle radius=9mm, angle eccentricity=1.35,
        pic text options={text=black,text opacity=1}] {angle = D--C--B};
    \pic [draw, fill=yellow!25, fill opacity=0.35, "$\beta$", angle radius=8mm, angle eccentricity=1.35,
        pic text options={text=black,text opacity=1}] {angle = C--B--O};
    
    \filldraw (O) circle (1pt) node[above right] {$O$};
    \node[below left] at (A) {$A$}; \node[below right] at (B) {$B$}; \node[above] at (C) {$C$};
\end{tikzpicture}
\end{center}

\textbf{3 Žingsnis (Priekampis):} Dabar pažiūrėkime į trikampį $AOC$. Taške $O$ jis turi išorinį priekampį $\angle AOD$. Pagal praeito handouto taisyklę, išorinis kampas lygus dviejų su juo negretimų vidaus kampų sumai:
$\angle AOD = \alpha + \alpha = \mathbf{2\alpha}$.
Lygiai taip pat, dešiniojo trikampio $BOC$ priekampis taške $O$ yra:
$\angle BOD = \beta + \beta = \mathbf{2\beta}$.

\begin{center}
\begin{tikzpicture}[scale=1.2]
    \coordinate (O) at (0,0);
    \draw[thick] (O) circle (2);
    \coordinate (A) at (210:2); \coordinate (B) at (330:2); \coordinate (C) at (70:2); \coordinate (D) at (250:2);
    \draw[thick, dashed] (C) -- (D); \draw[thick] (A) -- (C) -- (B); \draw[thick] (A) -- (O) -- (B); \draw[thick, blue] (O) -- (C);
    \foreach \p in {O,A,B,C,D} {\filldraw (\p) circle (1.2pt);}
    
    \pic [draw, fill=blue!20, fill opacity=0.35, "$\alpha$", angle radius=10mm, angle eccentricity=1.35,
        pic text options={text=black,text opacity=1}] {angle = A--C--D};
    \pic [draw, fill=blue!20, fill opacity=0.35, "$\alpha$", angle radius=8mm, angle eccentricity=1.35,
        pic text options={text=black,text opacity=1}] {angle = O--A--C};
    \pic [draw, fill=yellow!25, fill opacity=0.35, "$\beta$", angle radius=9mm, angle eccentricity=1.35,
        pic text options={text=black,text opacity=1}] {angle = D--C--B};
    \pic [draw, fill=yellow!25, fill opacity=0.35, "$\beta$", angle radius=8mm, angle eccentricity=1.35,
        pic text options={text=black,text opacity=1}] {angle = C--B--O};

    \pic [draw, fill=blue!20, fill opacity=0.35, "$2\alpha$", angle radius=8mm, angle eccentricity=1.35,
        pic text options={text=black,text opacity=1, xshift=-1mm}] {angle = A--O--D};
    \pic [draw, fill=yellow!25, fill opacity=0.35, "$2\beta$", angle radius=8mm, angle eccentricity=1.35,
        pic text options={text=black,text opacity=1, xshift=1mm}] {angle = D--O--B};
    
    \filldraw (O) circle (1pt);
    \node[above right] at (O) {$O$};
    \node[below left] at (A) {$A$}; \node[below right] at (B) {$B$}; \node[above] at (C) {$C$};
\end{tikzpicture}
\end{center}

\textbf{4 Žingsnis (Sujungimas):} 
Visas įbrėžtinis kampas $\angle ACB = \alpha + \beta$.
Visas centrinis kampas $\angle AOB = 2\alpha + 2\beta = 2(\alpha + \beta)$.
Vadinasi, centrinis kampas tikrai yra lygiai dvigubai didesnis. \hfill $\square$

\vspace{0.5cm}

% --- NAUJAS IŠVADŲ BLOKAS ---

Atkreipkite dėmesį, jog mums nereikia aklai kalti teorijos pradžioje minėtų savybių (2 ir 3). Šie du faktai \textbf{natūraliai išplaukia} iš ką tik įrodytos centrinio ir įbrėžtinio kampų savybės!

\begin{corollary}[Kampai, besiremiantys į tą patį lanką]
Jei nupiešime kelis skirtingus įbrėžtinius kampus, besiremiančius į lanką $AB$, jie visi turės \textbf{tą patį} centrinį kampą $\angle AOB$. Kadangi kiekvienas įbrėžtinis kampas pagal įrodytą taisyklę privalo būti lygiai perpus mažesnis už centrinį ($2\alpha \to \alpha$), tai visi įbrėžtiniai kampai privalo būti lygūs tarpusavyje!
\end{corollary}

\begin{center}
\begin{tikzpicture}[scale=1.3]
    \coordinate (O) at (0,0);
    \draw[thick] (O) circle (1.5);
    \coordinate (A) at (210:1.5);
    \coordinate (B) at (330:1.5);
    \coordinate (C1) at (50:1.5);
    \coordinate (C2) at (120:1.5);
    \foreach \p in {O,A,B,C1,C2} {\filldraw (\p) circle (1.2pt);}
    
    \draw[thick, dashed, red] (A) -- (O) -- (B);
    \draw[thick] (A) -- (C1) -- (B);
    \draw[thick] (A) -- (C2) -- (B);
    
    \pic [draw, fill=red!10, fill opacity=0.4, "$2\alpha$", angle radius=7mm, angle eccentricity=1.35,
        pic text options={text=black,text opacity=1}] {angle = A--O--B};
    \pic [draw, fill=blue!10, fill opacity=0.4, "$\alpha$", angle radius=6mm, pic text options={text=black,text opacity=1}] {angle = A--C1--B};
    \pic [draw, fill=blue!10, fill opacity=0.4, "$\alpha$", angle radius=6mm, pic text options={text=black,text opacity=1}] {angle = A--C2--B};
    \filldraw (O) circle (1pt) node[above, xshift=1mm] {$O$};
\end{tikzpicture}
\end{center}

\begin{corollary}[Kampas į skersmenį yra status]
Jei lankas yra lygus pusei apskritimo, tai šio lanko centrinis kampas yra tiesi linija – t.y., lygiai $180^\circ$. Pagal mūsų taisyklę, įbrėžtinis kampas privalo būti perpus mažesnis, todėl jis \textbf{visada} lygus $\frac{180^\circ}{2} = 90^\circ$!
\end{corollary}

\begin{center}
\begin{tikzpicture}[scale=1.3]
    \coordinate (O) at (0,0);
    \draw[thick] (O) circle (1.5);
    \coordinate (A) at (180:1.5);
    \coordinate (B) at (0:1.5);
    \coordinate (C) at (60:1.5);
    \foreach \p in {O,A,B,C} {\filldraw (\p) circle (1.2pt);}
    
    \draw[thick, red] (A) -- (B); % skersmuo
    \draw[thick] (A) -- (C) -- (B);
    
    \pic [draw, fill=red!10, fill opacity=0.4, "$180^\circ$", angle radius=3mm, angle eccentricity=1.35,
        pic text options={text=black,text opacity=1,yshift=1mm}] {angle = B--O--A};
    
    % Statusis kampas
    \coordinate (CdirA) at ($(C)!3mm!(A)$);
    \coordinate (CdirB) at ($(C)!3mm!(B)$);
    \coordinate (Crect) at ($(CdirA) + (CdirB) - (C)$);
    \draw[thick, fill=blue!10] (CdirA) -- (Crect) -- (CdirB) -- (C) -- cycle;
    
    \filldraw (O) circle (1pt) node[below] {$O$};
    \node[left] at (A) {$A$};
    \node[right] at (B) {$B$};
\end{tikzpicture}
\end{center}
\vspace{0.5cm}
% -----------------------------

\newpage
\begin{example}[Kampas tarp besikertančių stygų]
Taškai $A, B, C, D$ išsidėstę ant apskritimo. Stygos $AC$ ir $BD$ susikerta taške $E$. Duota, kad $\angle BAC = \alpha$, o $\angle ACD = \beta$. Įrodykite, kad kampas tarp dviejų besikertančių stygų $\angle BEC = \alpha + \beta$.
\end{example}

\begin{center}
\begin{tikzpicture}[scale=1.3]
    \draw[thick] (0,0) circle (2);
    \coordinate (A) at (150:2);
    \coordinate (B) at (30:2);
    \coordinate (C) at (320:2);
    \coordinate (D) at (210:2);
    \coordinate (E) at (intersection of A--C and B--D);
    \foreach \p in {A,B,C,D,E} {\filldraw (\p) circle (1.2pt);}

    \draw[thick] (A) -- (E) -- (C);
    \draw[thick] (B) -- (E) -- (D);
    \draw[thick, gray, dashed] (A) -- (B);
    \draw[thick, gray, dashed] (C) -- (D);

    \filldraw (E) circle (1.2pt);
    \pic[draw, fill=blue!20, fill opacity=0.35, "$\alpha$", angle radius=8mm, angle eccentricity=1.35,
        pic text options={text=black,text opacity=1}] {angle = C--A--B};
    \pic[draw, fill=yellow!25, fill opacity=0.35, "$\beta$", angle radius=8mm, angle eccentricity=1.35,
        pic text options={text=black,text opacity=1}] {angle = A--C--D};
    \node[above left] at (A) {$A$};
    \node[above right] at (B) {$B$};
    \node[below right] at (C) {$C$};
    \node[below left] at (D) {$D$};
    \node[above, yshift=1mm] at (E) {$E$};
\end{tikzpicture}
\end{center}

\textbf{Sprendimas:}\\
Sąlygoje neduotas kampas $\angle BEC$. Mes pasieksime jį pamažu.

\textbf{1 Žingsnis (Perkėlimas per lanką):} Kampai $\angle ACD$ ir $\angle ABD$ remiasi į tą patį lanką $AD$. Todėl pagal įbrėžtinių kampų lygybę:
\[
\angle ABD = \angle ACD = \mathbf{\beta}.
\]

\begin{center}
\begin{tikzpicture}[scale=1.3]
    \coordinate (O) at (0,0);
    \draw[thick] (O) circle (2);
    
    \coordinate (A) at (150:2);
    \coordinate (B) at (30:2);
    \coordinate (C) at (320:2);
    \coordinate (D) at (210:2);
    \coordinate (E) at (intersection of A--C and B--D);
    
    \draw[ultra thick, red] (150:2) arc (150:210:2); % lankas AD
    
    \draw[thick] (A) -- (C);
    \draw[thick] (B) -- (D);
    \draw[thick, gray, dashed] (C) -- (D);
    \draw[thick, gray, dashed] (A) -- (B);
    \foreach \p in {A,B,C,D,E} {\filldraw (\p) circle (1.2pt);}
    
    \pic [draw, fill=blue!20, fill opacity=0.35, "$\alpha$", angle radius=8mm, angle eccentricity=1.35,
        pic text options={text=black,text opacity=1}] {angle = C--A--B};
    \pic [draw, fill=yellow!25, fill opacity=0.35, "$\beta$", angle radius=8mm, angle eccentricity=1.35,
        pic text options={text=black,text opacity=1}] {angle = A--C--D};
    \pic [draw, fill=yellow!25, fill opacity=0.35, "$\beta$", angle radius=8mm, angle eccentricity=1.35,
        pic text options={text=black,text opacity=1, xshift=-2mm}] {angle = A--B--D};
    
    \node[above left] at (A) {$A$};
    \node[above right] at (B) {$B$};
    \node[below right] at (C) {$C$};
    \node[below left] at (D) {$D$};
    \node[above, yshift=1mm] at (E) {$E$};
\end{tikzpicture}
\end{center}

\textbf{2 Žingsnis (Priekampis ir pabaiga):} Susikoncentruokime į trikampį $\triangle ABE$.
Kadangi $E$ yra ant atkarpų $AC$ ir $BD$, gauname $\angle BAE = \angle BAC = \alpha$ ir $\angle ABE = \angle ABD = \beta$.
Kampas $\angle BEC$ yra trikampio $\triangle ABE$ išorinis kampas prie viršūnės $E$.
Pagal priekampio savybę:
\[
\angle BEC = \angle BAE + \angle ABE = \alpha + \beta.
\]
Įrodyta. \hfill $\square$

\newpage
\begin{example}[3. Pusiaukraštinė stačiajame trikampyje]
Trikampis $ABC$ yra statusis, $\angle C=90^\circ$. Taškas $M$ yra įžambinės $AB$ vidurys, todėl $AM=MB$. Įrodykite, kad $CM=AM=MB$.
\end{example}

\begin{center}
\begin{tikzpicture}[scale=1]
    \coordinate (A) at (-2,0); \coordinate (B) at (2,0);
    \coordinate (C) at (1.2,1.6); \coordinate (M) at (0,0);
    \draw[thick] (A) -- (C) -- (B) -- cycle;
    \draw[thick] (C) -- (M); \draw[thick] (A) -- (M) -- (B);
    \draw[thick, blue] ($(A)!0.5!(M)+(0,0.12)$) -- ($(A)!0.5!(M)+(0,-0.12)$);
    \draw[thick, blue] ($(M)!0.5!(B)+(0,0.12)$) -- ($(M)!0.5!(B)+(0,-0.12)$);
    \coordinate (CdirA) at ($(C)!4mm!(A)$);
    \coordinate (CdirB) at ($(C)!4mm!(B)$);
    \coordinate (Crect) at ($(CdirA) + (CdirB) - (C)$);
    \draw[thick, fill=red!20, fill opacity=0.35] (CdirA) -- (Crect) -- (CdirB) -- (C) -- cycle;
    \filldraw (A) circle (1.2pt); \filldraw (B) circle (1.2pt);
    \filldraw (C) circle (1.2pt); \filldraw (M) circle (1.2pt);
    \node[below] at (A) {$A$}; \node[below] at (B) {$B$};
    \node[above] at (C) {$C$}; \node[below] at (M) {$M$};
\end{tikzpicture}
\end{center}

\textbf{Sprendimas:} Šį uždavinį jau nagrinėjome ankstesniame handoute. Dabar jį išspręsime elegantiškiau, pasinaudodami apskritimo savybėmis.

\textbf{1 Žingsnis (Apskritimas):} Kadangi $\angle C=90^\circ$, pagal įbrėžtinio kampo, besiremiančio į skersmenį, savybę aplink trikampį $ABC$ galima apibrėžti apskritimą, kurio skersmuo yra $AB$.

\begin{center}
\begin{tikzpicture}[scale=1]
    \coordinate (M) at (0,0);
    \coordinate (A) at (-2,0);
    \coordinate (B) at (2,0);
    \coordinate (C) at (1.2,1.6);
    \draw[thick, dashed, gray] (M) circle (2);
    \draw[thick] (A)--(C)--(B)--cycle;
    \draw[thick] (M)--(C);
    \draw[thick, blue] ($(A)!0.5!(M)+(0,0.12)$)--($(A)!0.5!(M)+(0,-0.12)$);
    \draw[thick, blue] ($(M)!0.5!(B)+(0,0.12)$)--($(M)!0.5!(B)+(0,-0.12)$);
    \coordinate (CdirA) at ($(C)!4mm!(A)$);
    \coordinate (CdirB) at ($(C)!4mm!(B)$);
    \coordinate (Crect) at ($(CdirA)+(CdirB)-(C)$);
    \draw[thick, fill=red!20, fill opacity=0.35] (CdirA)--(Crect)--(CdirB)--(C)--cycle;
    \filldraw (M) circle (1.2pt); \filldraw (A) circle (1.2pt);
    \filldraw (B) circle (1.2pt); \filldraw (C) circle (1.2pt);
    \node[below] at (M) {$M$}; \node[left] at (A) {$A$};
    \node[right] at (B) {$B$}; \node[above] at (C) {$C$};
\end{tikzpicture}
\end{center}

\textbf{2 Žingsnis (Spinduliai):} Taškas $M$ yra skersmens $AB$ vidurys, todėl jis yra apibrėžto apskritimo centras. Vadinasi, $MA$, $MB$ ir $MC$ yra to paties apskritimo spinduliai.
\[
MA=MB=MC.
\]

\begin{center}
\begin{tikzpicture}[scale=1]
    \coordinate (M) at (0,0);
    \coordinate (A) at (-2,0);
    \coordinate (B) at (2,0);
    \coordinate (C) at (1.2,1.6);
    \draw[thick, dashed, gray] (M) circle (2);
    \draw[thick] (A)--(C)--(B)--cycle;
    \draw[thick] (M)--(C);
    \draw[thick, blue] ($(A)!0.5!(M)+(0,0.12)$)--($(A)!0.5!(M)+(0,-0.12)$);
    \draw[thick, blue] ($(M)!0.5!(B)+(0,0.12)$)--($(M)!0.5!(B)+(0,-0.12)$);
    \draw[thick, blue] ($(M)!0.5!(C)+(-0.08,0.06)$)--($(M)!0.5!(C)+(0.08,-0.06)$);
    \coordinate (CdirA) at ($(C)!4mm!(A)$);
    \coordinate (CdirB) at ($(C)!4mm!(B)$);
    \coordinate (Crect) at ($(CdirA)+(CdirB)-(C)$);
    \draw[thick, fill=red!20, fill opacity=0.35] (CdirA)--(Crect)--(CdirB)--(C)--cycle;
    \filldraw (M) circle (1.2pt); \filldraw (A) circle (1.2pt);
    \filldraw (B) circle (1.2pt); \filldraw (C) circle (1.2pt);
    \node[below] at (M) {$M$}; \node[left] at (A) {$A$};
    \node[right] at (B) {$B$}; \node[above] at (C) {$C$};
\end{tikzpicture}
\end{center}

\textbf{3 Žingsnis (Išvada):} Taigi pusiaukraštinė $CM$ lygi pusei įžambinės:
\[
CM=AM=MB.
\]
Įrodyta. \hfill $\square$

\vspace{0.5cm}

% ==========================================
% UŽDAVINIAI
% ==========================================
\newpage
\section{Uždaviniai savarankiškam sprendimui (12 iššūkių)}

\begin{enumerate}
    \renewcommand{\labelenumi}{\textbf{\theenumi.}}
    
    \item Apskritime nubrėžtas centrinis kampas $\angle AOB = 130^\circ$. Kokio dydžio yra įbrėžtinis kampas $\angle ACB$, jeigu taškas $C$ guli ant didesniojo apskritimo lanko?
    \item Apskritimo styga padalina apskritimą į du lankus. Iš vienos stygos pusės į ją remiasi įbrėžtinis kampas, lygus $70^\circ$. Kam lygus centrinis kampas, besiremiantis į šią stygą?
    \item Trikampis $ABC$ įbrėžtas į apskritimą. Žinoma, kad $\angle A = 50^\circ$, o $\angle B = 60^\circ$. Raskite visų trijų atitinkamų apskritimo lankų, į kuriuos remiasi šie kampai, dydžius (laipsniais).
    \item Taškai $A, B, C, D$ paeiliui išsidėstę ant apskritimo. Įrodykite, kad atsiradusių trikampių kampai dalijasi šiomis poromis: $\angle CAD = \angle CBD$ ir $\angle BCA = \angle BDA$.
    \item Taškai $A, B, C, D$ išsidėstę ant apskritimo, stygos $AC$ ir $BD$ susikerta taške $P$. Duota, kad $\angle BDC = 25^\circ$, o $\angle ACD = 45^\circ$. Apskaičiuokite kampą $\angle BPC$.
    \item Apskritime nubrėžtos dvi stygos $AB$ ir $CD$, kurios susikerta taške $K$. Žinoma, kad lankas $AC$ lygus $40^\circ$ (t.y. jo centrinis kampas $40^\circ$), o lankas $BD$ lygus $80^\circ$. Įrodykite, kad kampas tarp stygų $\angle AKC$ yra lygus lygiai pusei šių lankų sumos (t.y. $60^\circ$).
    \item Atkarpa $AB$ yra apskritimo skersmuo. Ant apskritimo pažymėtas taškas $C$. Iš taško $C$ nuleista aukštinė $CH$ į skersmenį $AB$. Žinoma, kad $\angle A = 35^\circ$. Raskite kampą $\angle BCH$.
    \item Trikampio $ABC$ kampai yra $60^\circ$, $70^\circ$ ir $50^\circ$. Šiame trikampyje nubrėžtos aukštinės $AD$ ir $BE$, kurios susikerta taške $H$. 
    \begin{itemize}
        \item[a)] Įrodykite, kad aplink keturkampį $CEHD$ galima apibrėžti apskritimą (kuris jo skersmuo?).
        \item[b)] Įrodykite, kad aplink keturkampį $ABDE$ irgi galima apibrėžti apskritimą.
    \end{itemize}
    \textit{(Šis uždavinys – tai įvadas į kitą mūsų modulį!)}
    \item Stačiojo trikampio $ABC$ ($\angle C = 90^\circ$) viduje pažymėtas įžambinės vidurio taškas $M$. Trikampio viduje nubrėžta aukštinė $CH$. Įrodykite, kad $\angle MCH = \angle A - \angle B$ (tariant, kad $\angle A > \angle B$).
    \item Trikampis $ABC$ įbrėžtas į apskritimą. Pusiaukampinė $AD$ kerta priešingą kraštinę $BC$ taške $D$, o patį apskritimą kerta taške $M$. Įrodykite, kad taškas $M$ yra lanko $BC$ vidurio taškas.
    \item Duotas smailusis trikampis $ABC$. Iš viršūnių $B$ ir $C$ nuleistos aukštinės $BE$ ir $CD$ kerta priešingas kraštines. Įrodykite, kad $\angle ADE = \angle ABC$.
    \item Apskritime nubrėžtos dvi lygiagrečios stygos $AB \parallel CD$. Įrodykite, kad tarp jų esantys lankai $AC$ ir $BD$ yra lygūs. (Vėlgi – naudokite tik atitinkamus įbrėžtinius kampus ir lygiagrečių tiesių savybę iš praeito handouto!).
\end{enumerate}

\end{document}